% Sección 6.2: Derivadas e Integrales de Funciones Exponenciales Naturales

\section*{SECCIÓN 6.2 \quad Ejercicios}

\subsection*{Ejercicios 1--2: Aproximaciones numéricas e interpretación analítica del límite definitorio de la base $e$.}

\begin{enumerate}
    \subimport{ejercicio_01/}{ejercicio.tex}
    \subimport{ejercicio_02/}{ejercicio.tex}
\end{enumerate}

\subsection*{Ejercicios 3--18: Diferenciación de funciones compuestas exponenciales naturales.}

\noindent \textit{Obtenga la derivada de las funciones dadas en los Ejercicios 3 a 18.}

\begin{enumerate}
    \setcounter{enumi}{2}
    \begin{minipage}{0.45\textwidth}
    \subimport{ejercicio_03/}{ejercicio.tex}
    \subimport{ejercicio_05/}{ejercicio.tex}
    \subimport{ejercicio_07/}{ejercicio.tex}
    \subimport{ejercicio_09/}{ejercicio.tex}
    \subimport{ejercicio_11/}{ejercicio.tex}
    \subimport{ejercicio_13/}{ejercicio.tex}
    \subimport{ejercicio_15/}{ejercicio.tex}
    \subimport{ejercicio_17/}{ejercicio.tex}
    \end{minipage}%
    \begin{minipage}{0.45\textwidth}
    \subimport{ejercicio_04/}{ejercicio.tex}
    \subimport{ejercicio_06/}{ejercicio.tex}
    \subimport{ejercicio_08/}{ejercicio.tex}
    \subimport{ejercicio_10/}{ejercicio.tex}
    \subimport{ejercicio_12/}{ejercicio.tex}
    \subimport{ejercicio_14/}{ejercicio.tex}
    \subimport{ejercicio_16/}{ejercicio.tex}
    \subimport{ejercicio_18/}{ejercicio.tex}
    \end{minipage}
\end{enumerate}

\subsection*{Ejercicios 19--20: Ecuación de la recta tangente en puntos coordenados dados.}

\noindent \textit{En los Ejercicios 19 y 20 encuentre la ecuación de la recta tangente a la curva dada en el punto indicado.}

\begin{enumerate}
    \setcounter{enumi}{18}
    \begin{minipage}{0.45\textwidth}
    \subimport{ejercicio_19/}{ejercicio.tex}
    \end{minipage}%
    \begin{minipage}{0.45\textwidth}
    \subimport{ejercicio_20/}{ejercicio.tex}
    \end{minipage}
\end{enumerate}

\subsection*{Ejercicios 21--32: Derivación implícita, ecuaciones diferenciales, derivadas de orden superior y modelado físico/numérico.}

\begin{enumerate}
    \setcounter{enumi}{20}
    \subimport{ejercicio_21/}{ejercicio.tex}
    \subimport{ejercicio_22/}{ejercicio.tex}
    \subimport{ejercicio_23/}{ejercicio.tex}
    \subimport{ejercicio_24/}{ejercicio.tex}
    \subimport{ejercicio_25/}{ejercicio.tex}
    \subimport{ejercicio_26/}{ejercicio.tex}
    \subimport{ejercicio_27/}{ejercicio.tex}
    \subimport{ejercicio_28/}{ejercicio.tex}
    \subimport{ejercicio_29/}{ejercicio.tex}
    \subimport{ejercicio_30/}{ejercicio.tex}
    \subimport{ejercicio_31/}{ejercicio.tex}
    \subimport{ejercicio_32/}{ejercicio.tex}
\end{enumerate}

\subsection*{Ejercicios 33--34: Gráficas mediante transformaciones de la curva natural.}

\noindent \textit{En los Ejercicios 33 y 34 utilice la gráfica de $y = e^x$ para trazar la gráfica de la función dada.}

\begin{enumerate}
    \setcounter{enumi}{32}
    \begin{minipage}{0.45\textwidth}
    \subimport{ejercicio_33/}{ejercicio.tex}
    \end{minipage}%
    \begin{minipage}{0.45\textwidth}
    \subimport{ejercicio_34/}{ejercicio.tex}
    \end{minipage}
\end{enumerate}

\subsection*{Ejercicios 35--40: Evaluación de límites infinitos, laterales y exponenciales trigonométricos.}

\noindent \textit{En los Ejercicios 35 a 40 encuentre el límite indicado.}

\begin{enumerate}
    \setcounter{enumi}{34}
    \begin{minipage}{0.45\textwidth}
    \subimport{ejercicio_35/}{ejercicio.tex}
    \subimport{ejercicio_37/}{ejercicio.tex}
    \subimport{ejercicio_39/}{ejercicio.tex}
    \end{minipage}%
    \begin{minipage}{0.45\textwidth}
    \subimport{ejercicio_36/}{ejercicio.tex}
    \subimport{ejercicio_38/}{ejercicio.tex}
    \subimport{ejercicio_40/}{ejercicio.tex}
    \end{minipage}
\end{enumerate}

\subsection*{Ejercicios 41--44: Extremos absolutos, intervalos de crecimiento, concavidad y puntos de inflexión.}

\begin{enumerate}
    \setcounter{enumi}{40}
    \subimport{ejercicio_41/}{ejercicio.tex}
    \subimport{ejercicio_42/}{ejercicio.tex}
\end{enumerate}

\noindent \textit{Para cada una de las funciones dadas en los Ejercicios 43 y 44 encuentre (a) los intervalos en donde es creciente o decreciente, (b) los intervalos de concavidad y (c) los puntos de inflexión.}

\begin{enumerate}
    \setcounter{enumi}{42}
    \begin{minipage}{0.45\textwidth}
    \subimport{ejercicio_43/}{ejercicio.tex}
    \end{minipage}%
    \begin{minipage}{0.45\textwidth}
    \subimport{ejercicio_44/}{ejercicio.tex}
    \end{minipage}
\end{enumerate}

\subsection*{Ejercicios 45--50: Análisis completo y trazo de curvas (Apartados A--H).}

\noindent \textit{En los Ejercicios 45 a 50 analice la curva dada conforme a los apartados A-H de la Sección 3.7.}

\begin{enumerate}
    \setcounter{enumi}{44}
    \begin{minipage}{0.45\textwidth}
    \subimport{ejercicio_45/}{ejercicio.tex}
    \subimport{ejercicio_47/}{ejercicio.tex}
    \subimport{ejercicio_49/}{ejercicio.tex}
    \end{minipage}%
    \begin{minipage}{0.45\textwidth}
    \subimport{ejercicio_46/}{ejercicio.tex}
    \subimport{ejercicio_48/}{ejercicio.tex}
    \subimport{ejercicio_50/}{ejercicio.tex}
    \end{minipage}
\end{enumerate}

\subsection*{Ejercicios 51--60: Integrales indefinidas de funciones exponenciales naturales.}

\noindent \textit{Evalúe las integrales dadas en los Ejercicios 51 a 60.}

\begin{enumerate}
    \setcounter{enumi}{50}
    \begin{minipage}{0.45\textwidth}
    \subimport{ejercicio_51/}{ejercicio.tex}
    \subimport{ejercicio_53/}{ejercicio.tex}
    \subimport{ejercicio_55/}{ejercicio.tex}
    \subimport{ejercicio_57/}{ejercicio.tex}
    \subimport{ejercicio_59/}{ejercicio.tex}
    \end{minipage}%
    \begin{minipage}{0.45\textwidth}
    \subimport{ejercicio_52/}{ejercicio.tex}
    \subimport{ejercicio_54/}{ejercicio.tex}
    \subimport{ejercicio_56/}{ejercicio.tex}
    \subimport{ejercicio_58/}{ejercicio.tex}
    \subimport{ejercicio_60/}{ejercicio.tex}
    \end{minipage}
\end{enumerate}

\subsection*{Ejercicios 61--67: Áreas, volúmenes de revolución, problemas inversos de valores iniciales y desigualdades fundamentales probadas por inducción.}

\begin{enumerate}
    \setcounter{enumi}{60}
    \subimport{ejercicio_61/}{ejercicio.tex}
    \subimport{ejercicio_62/}{ejercicio.tex}
    \subimport{ejercicio_63/}{ejercicio.tex}
    \subimport{ejercicio_64/}{ejercicio.tex}
    \subimport{ejercicio_65/}{ejercicio.tex}
    \subimport{ejercicio_66/}{ejercicio.tex}
    \subimport{ejercicio_67/}{ejercicio.tex}
\end{enumerate}
